\section{Analytic rigidity of zero resonance mechanisms}
\label{sec:analytic-resonance-rigidity}
The genericity theorem leaves open the stronger question whether every
nonconstant real-analytic observable has positive resonance variance.  Two
rigidity results substantially narrow any possible counterexample.  First,
pointwise vanishing of the Poisson bracket is impossible in the analytic
category.  Second, if the variance nevertheless vanishes, the cohomological
transfer inherits analyticity from the bracket.

\begin{theorem}[Analytic pointwise-bracket rigidity]
\label{thm:analytic-bracket-rigidity}
Let $S\in\mathrm{SL}(2,\mathbb Z)$ be hyperbolic and let
$Cg=g\circ S$.  If a real-analytic $g:\mathbb T^2\to\mathbb R$ satisfies
\[
 \{g,Cg\}\equiv0,
\]
then $g$ is constant.
\end{theorem}
\begin{proof}
Suppose $g$ is nonconstant and let
$R=\{x:dg_x\ne0\}$.  At least one first derivative of $g$ is a nonzero
real-analytic function, so its zero set has measure zero; hence $R$ is open
and has full measure.  On $R$ define the level-set tangent line
\[
 L(x)=\ker dg_x\in\mathbb{RP}^1.
\]
Whenever $x,Sx\in R$, the bracket identity says that the two nonzero
covectors $dg_x$ and $dg_{Sx}\circ S$ are linearly dependent.  Thus
\begin{equation}\label{eq:analytic-line-equivariance}
 L(Sx)=S L(x).
\end{equation}
The invariant full-measure set
$R_\infty=\bigcap_{n\in\mathbb Z}S^{-n}R$ contains almost every recurrent
point by Poincar\'e recurrence.  If $S^{n_j}x\to x$ for such a point, then
continuity gives $L(S^{n_j}x)\to L(x)$, while
\eqref{eq:analytic-line-equivariance} gives
$L(S^{n_j}x)=S^{n_j}L(x)$.

Let $E^s,E^u$ be the stable and unstable eigenlines.  The projective action
of a hyperbolic matrix has precisely these fixed points, and every line other
than $E^s$ converges under positive iteration to $E^u$.  Recurrence therefore
forces $L(x)\in\{E^s,E^u\}$ for almost every regular $x$.  By continuity this
holds everywhere on $R$.  On each connected component of $R$ the line field
is constant, because its target is the discrete two-point set
$\{E^s,E^u\}$.  On one nonempty component $U$ there is an eigenvector $e$
with
\[
 D_e g=0\quad\hbox{on }U.
\]
The left side is real analytic, hence vanishes on all of the connected torus.
An eigenline of a hyperbolic integral determinant-one matrix has irrational
slope: a rational eigenline would give a rational algebraic-integer
eigenvalue and therefore an integer unit $\pm1$.  The linear flow in the
direction $e$ is consequently dense.  Since $g$ is constant along that flow,
continuity makes $g$ constant on the torus, a contradiction.
\end{proof}

We now make the remaining cohomological alternative analytic as well.  Put
$A=S^{\mathsf T}$ for the frequency action.  After replacing $A$ by
$\operatorname{sign}(\operatorname{tr}A)A$, choose real eigen-covectors
$\xi,\eta$ with expanding eigenvalue $\lambda>1$ and define
\[
 N_A(k)=|\xi(k)|+|\eta(k)|.
\]
There are constants $0<c_A\le C_A$ such that
$c_A\|k\|_1\le N_A(k)\le C_A\|k\|_1$.

\begin{theorem}[Analytic orbit-sum coboundary regularity]
\label{thm:analytic-coboundary-regularity}
Let $h$ have zero mean and satisfy
\[
 |\widehat h(k)|\le M e^{-\rho N_A(k)}
 \qquad(k\ne0).
\]
If every oriented $A$-orbit has zero coefficient sum, then there is a transfer
$u$, unique modulo constants, such that
\[
 h=u-Cu,
\]
and
\begin{equation}\label{eq:analytic-transfer-decay}
 |\widehat u(k)|\le M' e^{-(\rho/2)N_A(k)}.
\end{equation}
In particular $u$ is real analytic whenever $h$ is real analytic.
\end{theorem}
\begin{proof}
Fix an orbit $k_j=A^jk_0$ and write
$X_j=|\xi(k_j)|$, $Y_j=|\eta(k_j)|$.  Then
$X_j=\lambda^jX_0$ and $Y_j=\lambda^{-j}Y_0$.
Choose the unique switching index $j_*$ with
$X_{j_*}\ge Y_{j_*}$ and $X_{j_*-1}<Y_{j_*-1}$.
For $j\ge j_*$ and $n\ge0$,
\[
 N_A(k_{j+n})\ge \frac{\lambda^n}{2}N_A(k_j),
\]
while for $j<j_*$,
\[
 N_A(k_{j-n})\ge \frac{\lambda^n}{2}N_A(k_j).
\]
The zero orbit sum allows the same transfer coefficient to be written as
\[
 u_j=\sum_{\ell\le j}h_\ell=-\sum_{\ell>j}h_\ell.
\]
Use the forward tail on the expanding side and the backward tail on the
contracting side.  The preceding inequalities give a uniformly convergent
geometric-superexponential bound by a constant times
$e^{-(\rho/2)N_A(k_j)}$.  This proves
\eqref{eq:analytic-transfer-decay} on every orbit.  The recurrence
$h_j=u_j-u_{j-1}$ is precisely the Fourier form of $h=u-Cu$.
If two decaying transfers existed, their difference would be $C$-invariant;
constant coefficients along every infinite nonzero frequency orbit are
incompatible with exponential decay, so only the zero frequency can remain.
\end{proof}

\begin{corollary}[Reduced all-analytic frontier]
\label{cor:analytic-resonance-frontier}
If a nonconstant real-analytic $g$ satisfies $\Gamma_A(g)=0$, then there is a
real-analytic $u$ such that
\begin{equation}\label{eq:analytic-nonlinear-cohomology}
 \frac{\{g,Cg\}}{4\pi}=u-Cu,
 \qquad \{g,Cg\}\not\equiv0.
\end{equation}
Thus an analytic zero-variance counterexample, if one exists, must solve the
nontrivial analytic nonlinear cohomological equation
\eqref{eq:analytic-nonlinear-cohomology}.  The compactly supported
Gevrey separated-support construction cannot realize this mechanism.
\end{corollary}

In ordinary Fourier coordinates, an estimate
$|\widehat h(k)|\le M e^{-r\|k\|_1}$ yields, for example,
\[
 |\widehat u(k)|\le M'
 e^{-r c_A\|k\|_1/(2C_A)}.
\]
These results stop short of the all-members analytic variance theorem. They
exclude the pointwise-zero mechanism and show that loss of transfer regularity
is not the remaining obstruction.
